\section{Supplementary evidence}
\subsection{Visual-copy audit}\label{app:audit}
The initial audit uses decoded RGB hashes, 64-bit perceptual hash Hamming distance at most six, and 16 by 16 RGB thumbnail mean absolute difference at most 18. Exact copies collapse; conflicting-label components are quarantined. The stronger audit adds quarter-turn/mirror hashes (same bit cutoff, aligned thumbnail difference at most 24), then proposes the nearest 15 frozen-feature neighbors within cosine distance 0.20.

ORB binary descriptors and a RANSAC homography verify proposed geometric links. Requirements are a 0.72 match-ratio cutoff, at least 16 filtered matches, 12 inliers, 65\% inlier fraction and 8\% convex-hull coverage in both images. This adds 948 transform and 1,797 geometric edges. There are no final conflicting components. All linked copies and initial hash components stay in one split. These checks reduce leakage but can still miss related photographs or join visually similar ones. Full criteria and preserved edges are in the commented audit script and data directory.

\subsection{Secondary diagnostics}\label{app:secondary}
\begin{table}[h]
\centering\small
\begin{tabular}{lrrr}
\toprule
Scores & NLL & ECE & Brier \\
\midrule
Unscaled & 0.859 & 0.043 & 0.470 \\
Scaled & 0.898 & 0.094 & 0.482 \\
\bottomrule
\end{tabular}
\caption{Unchanged contrast scores on test; lower is better. Scaling worsens NLL and ECE.}
\label{tab:calibration}
\end{table}

The original temperature is approximately 0.802. Validation is reused for regularization, temperature and review selection. The review search chooses the lowest threshold from 0.50 to 0.95 in steps of 0.05 attaining 90\% validation accuracy with at least 20 accepted images and 20\% coverage, with 0.80 fallback. Its chosen 0.70 threshold stays unchanged. At that threshold confidence alone covers 62.2\% with 77.6\% accepted accuracy; adding agreement covers 58.9\% with 77.7\%. This modest difference is descriptive, not a guaranteed improvement.

\begin{table}[h]
\centering\small
\begin{tabular}{lrrrr}
\toprule
Head & Clean & Dark & Blur & Mirror \\
\midrule
CNN global & 63.0 & 62.8 & 48.0 & 62.9 \\
Contrast & 65.5 & 67.1 & 49.7 & 65.2 \\
\bottomrule
\end{tabular}
\caption{Preserved descriptive stress-test accuracy in percent; no retraining or parameter choice.}
\label{tab:stress}
\end{table}

Dark multiplies source brightness by 0.6; blur uses Gaussian radius two at source resolution; mirror reflects horizontally while retaining the four broad labels. Both CNN views are recomputed. These conditions do not represent a new match or camera, and blur severity varies with source resolution. Original predictions, stress evidence and CNN-only timings remain under \texttt{results/images/}; preliminary hash-only evidence remains under \texttt{results/images\_initial\_hash/}.

\subsection{Exploratory video pilot}\label{app:video}
The separate Kaggle CRICKET\_SHOTS archive~\cite{dataset} contains 562 clips in eight labels and 25 normalized source-title groups. Its test set has 116 clips from five groups. Nine uniformly sampled frames feed the frozen encoder. Three ordered phase means, late-minus-early difference and temporal standard deviation form 6,400 features. This sparse temporal summary is motivated by segment-based representations~\cite{wang}, without reproducing TSN training.

The full-frame pilot obtains 18.1\% accuracy. A fixed upper-center crop, introduced after pilot test feedback, reaches 30.2\% accuracy and 21.3\% macro F1. The training-majority baseline also reaches 30.2\% accuracy (5.8\% macro F1). Equal majority accuracy is not a successful general video recognizer. Defence has only three test clips; validation has one unorthodox clip. The reused holdout makes this exploratory. Original and revised predictions remain available; different data and labels prevent a direct image-versus-video comparison.

\subsection{Code, data and repositories}
The complete report appends every current custom executable source and configuration, with a SHA-256 inventory. Public libraries, pretrained weights and the unmodified CVPR style are not original code. The project includes dependency files, manifests, numeric heads, saved predictions, tests and reproduction commands. The separate course mapping and viva guide explain lecture references, advanced audit tools and one measured image trace.

Project repository: \url{https://github.com/hrushikeshroop/cricket-shot-recognition-ee604}. Codex assisted with implementation, auditing, experiments and drafting. Members must review the implementation, evidence and limitations under course rules.

Data: \href{https://www.kaggle.com/datasets/aneesh10/cricket-shot-dataset}{Cricket Shots Dataset v2} and \href{https://www.kaggle.com/datasets/amittalmale/cricket-shots}{CRICKET\_SHOTS v1}. Pretrained weights: \href{https://github.com/onnx/models}{ONNX Model Zoo}. Libraries: \href{https://github.com/scikit-learn/scikit-learn}{scikit-learn}, \href{https://github.com/microsoft/onnxruntime}{ONNX Runtime}, \href{https://github.com/opencv/opencv}{OpenCV}, \href{https://github.com/pallets/flask}{Flask}, and \href{https://github.com/numpy/numpy}{NumPy}. Complete dependency links and attribution accompany \texttt{docs/sources.md}. Our NumPy HOG follows the cited method and supplied lectures; no external cricket classifier code was imported.
